\documentclass[aspectratio=169,8pt]{beamer}

\usetheme{default}
\usefonttheme{professionalfonts}
\setbeamertemplate{navigation symbols}{}
\setbeamertemplate{footline}{\hfill\scriptsize\insertframenumber/\inserttotalframenumber\hspace{2mm}\vspace{1mm}}

\usepackage{amsmath,amssymb,bm}
\usepackage{mathtools}
\usepackage{xcolor}
\usepackage[most]{tcolorbox}

\definecolor{deepblue}{RGB}{0,70,150}
\definecolor{softblue}{RGB}{232,242,255}
\definecolor{softgreen}{RGB}{235,250,238}
\definecolor{softorange}{RGB}{255,245,230}
\definecolor{softpurple}{RGB}{245,238,255}
\definecolor{darkred}{RGB}{160,35,25}
\definecolor{darkgreen}{RGB}{20,120,60}
\definecolor{darkpurple}{RGB}{95,45,150}

\setbeamerfont{frametitle}{series=\bfseries,size=\large}
\setbeamercolor{frametitle}{fg=black}
\setbeamertemplate{itemize item}{\color{deepblue}$\blacktriangleright$}
\setlength{\abovedisplayskip}{3pt}
\setlength{\belowdisplayskip}{3pt}

\newtcolorbox{bluebox}[1][]{colback=softblue,colframe=deepblue,boxrule=0.8pt,arc=1.5mm,left=1.5mm,right=1.5mm,top=0.8mm,bottom=0.8mm,#1}
\newtcolorbox{greenbox}[1][]{colback=softgreen,colframe=darkgreen,boxrule=0.8pt,arc=1.5mm,left=1.5mm,right=1.5mm,top=0.8mm,bottom=0.8mm,#1}
\newtcolorbox{orangebox}[1][]{colback=softorange,colframe=darkred,boxrule=0.8pt,arc=1.5mm,left=1.5mm,right=1.5mm,top=0.8mm,bottom=0.8mm,#1}
\newtcolorbox{purplebox}[1][]{colback=softpurple,colframe=darkpurple,boxrule=0.8pt,arc=1.5mm,left=1.5mm,right=1.5mm,top=0.8mm,bottom=0.8mm,#1}

\newcommand{\Cj}{C_j}
\newcommand{\Jj}{\bm{J}_j}
\newcommand{\xvec}{\bm{x}}
\newcommand{\nvec}{\bm{n}}
\newcommand{\kB}{k_{\mathrm B}}
\newcommand{\OmegaD}{\Omega}

\begin{document}

\begin{frame}{Module 3: Defect evolution after radiation deposition}
\scriptsize
\vspace{-1mm}
\begin{bluebox}
\centering\small
\textbf{Goal:} convert the Module 1 radiation damage map into time-dependent defect concentrations that later perturb space charge, temperature, recombination, and mobility.
\end{bluebox}

\vspace{1mm}
\begin{columns}[T,totalwidth=\textwidth]
\begin{column}{0.49\textwidth}
\textbf{Defect conservation law}
\begin{equation*}
\boxed{\frac{\partial C_j}{\partial t}+\nabla\cdot\Jj=S_j+R_j}
\end{equation*}
\scriptsize
\begin{itemize}\setlength\itemsep{0.2mm}
\item $C_j(x,y,t)$: concentration of defect species $j$.
\item $\Jj$: flux of species $j$ through the silicon lattice.
\item $S_j$: generation source from radiation-induced damage.
\item $R_j$: net reaction term: recombination, clustering, dissociation, annealing.
\end{itemize}

\vspace{1mm}
\textbf{Diffusive migration}
\begin{equation*}
\boxed{\Jj=-D_j\nabla C_j}
\end{equation*}
\scriptsize
Flux points from high defect concentration to low defect concentration.
\end{column}

\begin{column}{0.49\textwidth}
\textbf{Core diffusion-reaction equation}
\begin{equation*}
\boxed{\frac{\partial C_j}{\partial t}
=\nabla\cdot\left(D_j\nabla C_j\right)+S_j+R_j}
\end{equation*}

\textbf{2D project form}
\begin{equation*}
\boxed{
\frac{\partial C_j}{\partial t}
=\frac{\partial}{\partial x}\!\left(D_j\frac{\partial C_j}{\partial x}\right)
+\frac{\partial}{\partial y}\!\left(D_j\frac{\partial C_j}{\partial y}\right)
+S_j+R_j}
\end{equation*}
\scriptsize
\begin{itemize}\setlength\itemsep{0.2mm}
\item Divergence form is needed when $D_j=D_j(x,y,T,C_1,\ldots)$ varies spatially.
\item For constant $D_j$ and no source after irradiation:
\end{itemize}
\begin{equation*}
\frac{\partial C_j}{\partial t}=D_j\nabla^2 C_j+R_j.
\end{equation*}
\end{column}
\end{columns}

\vspace{0.5mm}
\begin{orangebox}
\centering\scriptsize
\textbf{Module handoff:} Module 3 outputs $C_j(x,y,t)$, the evolving defect map used by electrostatics, thermal transport, and carrier transport.
\end{orangebox}
\end{frame}

\begin{frame}{Module 3: Recombination, annealing, and temperature-dependent kinetics}
\scriptsize
\vspace{-1mm}
\begin{bluebox}
\centering\small
The reaction term $R_j$ is where lattice damage becomes material evolution: defects can annihilate, transform, cluster, or anneal.
\end{bluebox}

\vspace{1mm}
\begin{columns}[T,totalwidth=\textwidth]
\begin{column}{0.49\textwidth}
\textbf{Reduced single-species annealing}
\begin{equation*}
\boxed{R=-k_{\mathrm{ann}}C}
\end{equation*}
\begin{equation*}
\boxed{\frac{\partial C}{\partial t}=\nabla\cdot(D\nabla C)-k_{\mathrm{ann}}C}
\end{equation*}
\scriptsize
If diffusion is absent:
\begin{equation*}
\boxed{C(t)=C(0)e^{-k_{\mathrm{ann}}t}},\qquad
\boxed{\tau_{\mathrm{ann}}=k_{\mathrm{ann}}^{-1}}.
\end{equation*}

\vspace{1mm}
\textbf{Vacancy-interstitial recombination}
\begin{equation*}
V+I\longrightarrow \varnothing
\end{equation*}
\begin{equation*}
\boxed{R_V=-k_{VI}C_VC_I,\qquad R_I=-k_{VI}C_VC_I}
\end{equation*}
\end{column}

\begin{column}{0.49\textwidth}
\textbf{Arrhenius migration and annealing rates}
\begin{equation*}
\boxed{D_j(T)=D_{0,j}\exp\!\left[-\frac{E_{D,j}}{\kB T}\right]}
\end{equation*}
\begin{equation*}
\boxed{k_{\mathrm{ann},j}(T)=k_{0,j}\exp\!\left[-\frac{E_{\mathrm{ann},j}}{\kB T}\right]}
\end{equation*}
\scriptsize
\begin{itemize}\setlength\itemsep{0.2mm}
\item $E_{D,j}$: migration barrier for defect species $j$.
\item $E_{\mathrm{ann},j}$: barrier for removal/transformation channel.
\item Higher $T$ from Module 4 can sharply accelerate both diffusion and annealing.
\end{itemize}

\vspace{1mm}
\textbf{Transport timescale}
\begin{equation*}
\boxed{\tau_{D,j}\sim \frac{L^2}{D_j(T)}}
\end{equation*}
\small
The relevant competition is not a universal order; it is set by $D_j(T)$, $k_{\mathrm{ann},j}(T)$, concentrations, and geometry.
\end{column}
\end{columns}

\vspace{0.5mm}
\begin{purplebox}
\centering\scriptsize
\textbf{Key physics:} the same radiation damage map can either persist, spread, recombine, or anneal depending on temperature, defect species, and reaction pathways.
\end{purplebox}
\end{frame}

\begin{frame}{Module 3: FEM form, verification limits, and coupling role}
\scriptsize
\vspace{-1mm}
\begin{bluebox}
\centering\small
For implementation, Module 3 is a 2D transient diffusion-reaction FEM problem with source maps from Module 1 and temperature-dependent coefficients from Module 4.
\end{bluebox}

\vspace{1mm}
\begin{columns}[T,totalwidth=\textwidth]
\begin{column}{0.50\textwidth}
\textbf{Weak form for a defect species}
\begin{equation*}
\boxed{
\int_{\OmegaD} v\frac{\partial C_j}{\partial t}\,d\Omega
+\int_{\OmegaD} D_j\nabla v\cdot\nabla C_j\,d\Omega
=\int_{\OmegaD} v(S_j+R_j)\,d\Omega
}
\end{equation*}
\scriptsize
for all test functions $v$. A common default boundary condition is zero defect flux:
\begin{equation*}
\boxed{\nvec\cdot\Jj=0\quad\Longleftrightarrow\quad \nvec\cdot D_j\nabla C_j=0.}
\end{equation*}

\textbf{Initial condition from Module 1}
\begin{equation*}
\boxed{C_j(x,y,0)=C_{j,0}(x,y)\quad \text{or}\quad S_j=S_j(x,y,t).}
\end{equation*}
\end{column}

\begin{column}{0.48\textwidth}
\textbf{Verification cases}
\begin{greenbox}
\scriptsize
\textbf{Pure annealing:}
\begin{equation*}
D=0,\quad \frac{\partial C}{\partial t}=-k_{\mathrm{ann}}C
\quad\Rightarrow\quad C=C_0e^{-k_{\mathrm{ann}}t}.
\end{equation*}
\end{greenbox}
\vspace{0.7mm}
\begin{greenbox}
\scriptsize
\textbf{Pure diffusion:}
\begin{equation*}
R=0,\quad \frac{d}{dt}\int_{\OmegaD}C\,d\Omega=0
\quad \text{for zero-flux boundaries.}
\end{equation*}
\end{greenbox}
\vspace{0.7mm}
\begin{orangebox}
\scriptsize
\textbf{Coupled-module meaning:} $C_j$ modifies
\begin{equation*}
\rho_{\mathrm{def}}=q\sum_j z_jC_j,
\qquad
\mu_n,\mu_p,
\qquad
R_{\mathrm{SRH}},
\qquad
\dot q_{\mathrm{th}}.
\end{equation*}
Thus defect evolution changes electric fields, current flow, recombination, and heating.
\end{orangebox}
\end{column}
\end{columns}

\vspace{0.5mm}
\begin{purplebox}
\centering\scriptsize
\textbf{Summary:} Module 3 is the kinetic bridge between ``where radiation deposits energy'' and ``where the device remains electrically and thermally damaged.''
\end{purplebox}
\end{frame}

\end{document}
